\section*{Bab \ref{ch:b2:orbital-forcing} --- Gaya Paksa Orbital dan Zaman Es}

\begin{solution}{exo:b2:orbital-forcing:1}
(a) $1.0167/0.9833 = 1.0340$: aphelion \qty{3.4}{\percent} lebih jauh daripada perihelion.
(b) $a(1 + e) - a(1 - e) = 2ea = 0.0334\,a$, \qty{3.34}{\percent} dari $a$, dekat dengan 3.4 persen
yang biasa dikutip. (c) Sinar Matahari melemah sebanding dengan kebalikan kuadrat jarak: $1.0340^2 = 1.069$,
sekitar \qty{7}{\percent} lebih banyak sinar Matahari di perihelion.
\end{solution}

\begin{solution}{exo:b2:orbital-forcing:2}
(a) Kini $0.017236 \times \sin\ang{101.37} = 0.017236 \times 0.9804 = 0.01690$; \num{10000} tahun yang lalu
$0.019971 \times \sin\ang{293.86} = 0.019971 \times (-0.9145) = -0.01826$; \num{6000} tahun yang lalu
$0.019249 \times \sin\ang{359.99} \approx 0$, yaitu nilai-nilai dalam tabel. (b) Bujur Matahari pada saat
perihelion adalah $\varpi + \ang{180}$: \ang{281.37} masa kini, \ang{473.86} $-$ \ang{360} $=$ \ang{113.86}
\num{10000} tahun yang lalu, \ang{179.99} \num{6000} tahun yang lalu. (c) \num{10000} tahun yang lalu perihelion
jatuh \ang{24} setelah solstis Juni, yang terdekat dengannya di antara ketiganya, pada musim panas utara;
\num{6000} tahun yang lalu perihelion jatuh pada ekuinoks September, tempat \omterm{def:b2:orbital-forcing:elements}{presesi iklim} bernilai nol.
\end{solution}

\begin{solution}{exo:b2:orbital-forcing:3}
(a) Frekuensinya dijumlahkan: $1/25.77 + 1/112 = 0.03880 + 0.00893 = \qty{0.04773}{ka^{-1}}$, periode
sebesar $1/0.04773 = 21.0$ ribu tahun. (b) Hasilnya jatuh di dalam rentang periode utama,
19.0 sampai 23.7 ribu tahun. Satu angka tunggal tidak dapat mereproduksi semuanya: perputaran orbit
itu sendiri merupakan jumlah beberapa gerak dengan frekuensi berbeda, yang masing-masing bergabung dengan
presesi sumbu untuk memberikan satu periode \omterm{def:b2:orbital-forcing:elements}{presesi iklim}; rata-ratanya adalah
\num{23000} tahun yang biasa dikutip.
\end{solution}

\begin{solution}{exo:b2:orbital-forcing:4}
$\sqrt{1 - 0.0167^2} = 0.99986$. Untuk \ang{22.1}: $1360 \times 0.37622/(\pi \times 0.99986) =
\qty{162.9}{W.m^{-2}}$; untuk \ang{24.5}: $1360 \times 0.41469/(\pi \times 0.99986) =
\qty{179.5}{W.m^{-2}}$. Perbandingannya $\sin\ang{24.5}/\sin\ang{22.1} = 1.102$: rata-rata tahunan di
kutub bertambah \qty{10.2}{\percent}, sekitar \qty{16.7}{W.m^{-2}}, di antara kedua ekstrem itu.
\end{solution}

\begin{solution}{exo:b2:orbital-forcing:5}
(a) Pada solstis Juni berlaku $(a/r)^2 = \bigl((1 - e\sin\varpi)/(1 - e^2)\bigr)^2$. Kini $e\sin\varpi
= 0.016898$ dan $(0.983102/0.999703)^2 = 0.9671$, dibandingkan dengan $1 - 0.033796 = 0.9662$;
\num{116000} tahun yang lalu $e\sin\varpi = 0.043939$ dan $(0.956061/0.998065)^2 = 0.9176$, dibandingkan dengan $1 -
0.087878 = 0.9121$. (b) Pada solstis Desember $(a/r)^2 = \bigl((1 + e\sin\varpi)/(1 -
e^2)\bigr)^2$: kini $1.0347$ dibandingkan dengan $1.0338$, \num{116000} tahun yang lalu $1.0940$ dibandingkan dengan $1.0879$.
(c) Suku-suku yang diabaikan berorde $e^2$: $(1 - x)^2/(1 - e^2)^2 \approx 1 - 2x + x^2 + 2e^2$
dengan $x = e\sin\varpi$, di sini sekitar $3e^2$. Suku-suku itu bernilai $3 \times 0.0172^2 = 0.0009$ masa kini dan $3
\times 0.0440^2 = 0.0058$ \num{116000} tahun yang lalu, dekat dengan selisih yang ditemukan, $0.0009$ dan
$0.0055$: pendekatan itu enam kali lebih buruk ketika eksentrisitasnya dua setengah kali
lebih besar.
\end{solution}

\begin{solution}{exo:b2:orbital-forcing:6}
Pada solstis Juni $\delta = \varepsilon = \ang{23.446}$. $\cos h_0 = -\tan\ang{65}\tan\ang{23.446}
= -2.14451 \times 0.43369 = -0.93006$, sehingga $h_0 = \ang{158.44} = \qty{2.7654}{rad}$. Siang hari berlangsung
$2 \times 158.44/15 = 21.13$~jam (Bumi berputar \ang{15} per jam). Dengan $\sin\ang{65} = 0.90631$,
$\cos\ang{65} = 0.42262$, $\sin\delta = 0.39788$, $\cos\delta = 0.91744$, dan $\sin h_0 = 0.36742$,
bagian dalam kurung adalah $2.7654 \times 0.90631 \times 0.39788 + 0.42262 \times 0.91744 \times 0.36742 = 0.99720
+ 0.14246 = 1.13966$. Pada jarak rata-rata $\bar Q = (1360/\pi) \times 1.13966 = 432.90 \times 1.13966
= \qty{493.4}{W.m^{-2}}$; dengan faktor jarak dari \cref{exo:b2:orbital-forcing:5}, $493.4 \times
0.9671 = \qty{477.1}{W.m^{-2}}$. Sebagai pemeriksaan, langkah-langkah yang sama pada \ang{60} LU memberikan \qty{475.96}{W.m^{-2}}, sedangkan tabel-tabel itu memberikan \num{475.95}.
\end{solution}

\begin{solution}{exo:b2:orbital-forcing:7}
(a) Ekuinoks Maret adalah $\lambda = 0$, ekuinoks September $\lambda = \pi$. Dengan $\sqrt{1 - e^2}
\approx 1$, durasinya $\frac{T}{2\pi}\int_0^{\pi}\bigl(1 + 2e\cos(\lambda -
\varpi)\bigr)\,\mathrm d\lambda = \frac{T}{2\pi}\bigl(\pi + 2e[\sin(\lambda - \varpi)]_0^{\pi}\bigr)
= \frac{T}{2\pi}(\pi + 4e\sin\varpi) = \frac T2 + \frac{2T}{\pi}e\sin\varpi$. (b) Setengah tahun yang lain
lebih pendek dengan jumlah yang sama, sehingga selisihnya $(4T/\pi)\,e\sin\varpi = (4 \times
365.2422/\pi) \times 0.016898 = 465.04 \times 0.016898 = 7.86$~hari. (c) Musim semi dan musim panas bersama-sama
melebihi musim gugur dan musim dingin sebesar $4.5 + 3 = 7.5$~hari; kedua selisih yang dikutip itu dibulatkan, dan
kesesuaiannya berada di dalam pembulatan itu.
\end{solution}

\begin{solution}{exo:b2:orbital-forcing:8}
(a) Untuk $e = 0.0046$: $\sqrt{1 - e^2} = 0.999989$ dan $340/0.999989 = \qty{340.004}{W.m^{-2}}$. Untuk
$e = 0.0565$: $\sqrt{1 - 0.003192} = 0.998403$ dan $340/0.998403 = \qty{340.544}{W.m^{-2}}$. (b)
Selisihnya \qty{0.54}{W.m^{-2}}, \qty{0.16}{\percent}. (c) Insolasi pertengahan Juli pada \ang{65} LU
berayun sebesar $497 - 388 = \qty{109}{W.m^{-2}}$, sekitar \qty{25}{\percent} dari rata-ratanya dan dua ratus
kali perubahan rata-rata tahunan global: orbit memengaruhi iklim dengan memindahkan sinar Matahari antara
musim dan lintang, bukan dengan mengubah totalnya.
\end{solution}

\begin{solution}{exo:b2:orbital-forcing:9}
Dengan $\omega = 2\pi/P$ dan $\tau = \qty{10}{ka}$: untuk $P = 23$, $\omega\tau = 2.732$, amplitudo
relatif $1/\sqrt{1 + 7.464} = 0.344$ dan tunda $\arctan(2.732)/\omega = 1.2199/0.2732 = 4.47$~ka;
untuk $P = 41$, $\omega\tau = 1.532$, amplitudo $0.546$, tunda $0.9926/0.1532 = 6.48$~ka; untuk $P =
100$, $\omega\tau = 0.628$, amplitudo $0.847$, tunda $0.5610/0.0628 = 8.93$~ka. Lembaran es
mengunggulkan gaya paksa yang lambat: lembaran es mengikuti \qty{85}{\percent} siklus \num{100000} tahun tetapi hanya
sepertiga siklus presesi, karena waktunya terlalu sedikit untuk tumbuh atau menyusut sebelum gaya paksa yang cepat
berbalik arah. Lembaran es bertindak sebagai tapis lolos-rendah.
\end{solution}

\begin{solution}{exo:b2:orbital-forcing:10}
(a) Persamaan itu linear: jika $V_1$ menanggapi $A_1\sin\omega_1t$ dan $V_2$ menanggapi
$A_2\sin\omega_2t$, $V_1 + V_2$ menanggapi jumlahnya. Menurut \cref{prop:b2:orbital-forcing:lag} $V_1$ dan
$V_2$ adalah sinusoid dengan frekuensi $\omega_1$ dan $\omega_2$ (setelah transiennya lenyap), sehingga
tanggapan itu tidak memuat frekuensi lain. (b) Dengan $V \approx B_1\sin(\omega_1t - \phi_1) +
B_2\sin(\omega_2t - \phi_2)$,
\[
\begin{aligned}
  V^2 &= \tfrac12(B_1^2 + B_2^2) - \tfrac12B_1^2\cos(2\omega_1t - 2\phi_1) -
  \tfrac12B_2^2\cos(2\omega_2t - 2\phi_2)\\
  &\quad + B_1B_2\bigl[\cos\bigl((\omega_1 - \omega_2)t - \phi_1 + \phi_2\bigr) -
  \cos\bigl((\omega_1 + \omega_2)t - \phi_1 - \phi_2\bigr)\bigr],
\end{aligned}
\]
sehingga suku $kV^2$ memaksa persamaan itu pada frekuensi selisih dan frekuensi jumlah, serta pada frekuensi
ganda. (c) $1/(1/19.0 - 1/23.7) = 1/(0.05263 - 0.04219) = 95.8$ ribu tahun; pasangan 19.0
dan 22.4 memberikan 125 ribu tahun. Frekuensi selisih suku-suku presesi jatuh dalam pita
eksentrisitas, sebagaimana seharusnya, karena \omterm{def:b2:orbital-forcing:elements}{presesi iklim} adalah $e\sin\varpi$. Karena itu lembaran es yang tak linear
dapat menarik irama \num{100000} tahun dari presesi; lembaran es yang linear
tidak dapat. Siklus \num{100000} tahun merupakan bukti perilaku tak linear, meskipun perhitungan ini
saja tidak menunjukkan perilaku yang mana.
\end{solution}

\begin{solution}{exo:b2:orbital-forcing:11}
(a) Seperti dalam bukti \cref{prop:b2:orbital-forcing:annual}, $(a/r)^2\,\mathrm dt =
T\,\mathrm d\lambda/(2\pi\sqrt{1 - e^2})$, dan $\sin\delta = \sin\varepsilon\sin\lambda$ positif
untuk $0 < \lambda < \pi$. Jadi $E = \frac{S_0T\sin\varepsilon}{2\pi\sqrt{1 - e^2}}\int_0^{\pi}\sin\lambda
\,\mathrm d\lambda = S_0T\sin\varepsilon/(\pi\sqrt{1 - e^2})$. (b) $T = 365.25 \times 86400 =
\qty{3.1558e7}{s}$, $S_0T/\pi = \qty{1.3661e10}{J.m^{-2}}$; untuk \ang{22.1}, $E = 1.3661 \times 10^{10}
\times 0.37622/0.99986 = \qty{5.14}{GJ.m^{-2}}$; untuk \ang{24.5}, $\qty{5.67}{GJ.m^{-2}}$,
\qty{10}{\percent} lebih banyak. (c) Presesi menentukan bagaimana energi yang sama tersebar dalam waktu: musim panas
di perihelion lebih singkat dan lebih intens, musim panas di aphelion lebih panjang dan lebih lemah. Pencairan baru dimulai
di atas \qty{0}{\celsius}, sehingga musim panas yang singkat dan intens dapat mencairkan lebih banyak daripada musim panas yang panjang dan sejuk
yang menerima total yang sama.
\end{solution}

\begin{solution}{exo:b2:orbital-forcing:12}
(a) $3.93\log_2(279.4/189.7) = 3.93\log_2 1.4729 = 3.93 \times 0.5586 = \qty{2.20}{W.m^{-2}}$. (b)
Kenaikan terkecil adalah dari 194 ke \qty{266}{ppm}: $3.93\log_2 1.3711 = 3.93 \times 0.4554 =
\qty{1.79}{W.m^{-2}}$; yang terbesar dari 188 ke \qty{282}{ppm}: $3.93\log_2 1.5 = 3.93 \times 0.5850 =
\qty{2.30}{W.m^{-2}}$. (c) Gaya radiatif karbon dioksida empat kali perubahan terbesar
insolasi tahunan global, \qty{0.54}{W.m^{-2}}, dan bekerja di mana-mana dan sepanjang tahun; insolasi musim panas
pada \ang{65} LU berayun sekitar \qty{109}{W.m^{-2}}, tetapi pada satu lintang dan satu musim.
Orbit adalah pemicunya, melalui musim panas lembaran es utara; karbon dioksida, bersama
\omterm{def:b2:precambrian-earth:ice-albedo}{umpan balik es--albedo}, adalah penguat yang menyebarkan perubahan itu ke seluruh bola Bumi.
\end{solution}

\begin{solution}{pb:b2:orbital-forcing:1}
\textbf{1.} Dulu $0.0440 \times \sin\ang{92.7} = 0.0440 \times 0.99889 = 0.04395$; kini $0.0172 \times
\sin\ang{101.4} = 0.0172 \times 0.98027 = 0.01686$, sekitar 2.6 kali lebih kecil.
\textbf{2.} Dulu $1 - 0.0440 = 0.956$ dan $1 + 0.0440 = 1.044$; kini $0.9828$ dan $1.0172$.
\textbf{3.} Pada $\lambda = \varpi + \ang{180}$: dulu \ang{272.7}, kini \ang{281.4}, keduanya dekat
solstis Desember ($\lambda = \ang{270}$); dulu Bumi melewati perihelion kurang dari tiga hari
setelahnya, dan aphelion pada solstis Juni.
\textbf{4.} Dulu $(1.044/0.956)^2 = 1.0921^2 = 1.193$, kini $(1.0172/0.9828)^2 = 1.0350^2 = 1.071$:
\qty{19}{\percent} lebih banyak sinar Matahari di perihelion dulu, \qty{7}{\percent} kini.
\textbf{5.} Lebih kecil, \ang{22.52} dibandingkan dengan \ang{23.45}: Matahari musim panas berdiri lebih rendah di lintang tinggi,
dan lintang tinggi menerima lebih sedikit sinar Matahari pada musim panas dan sepanjang tahun.
\textbf{6.} Dulu $340/\sqrt{1 - 0.001936} = 340/0.999032 = \qty{340.33}{W.m^{-2}}$; kini $340/0.999852
= \qty{340.05}{W.m^{-2}}$. Selisihnya, \qty{0.28}{W.m^{-2}} atau kurang dari satu per seribu,
tidak berarti.
\textbf{7.} Pada $\lambda = \ang{90}$, $\cos(\ang{90} - \varpi) = \sin\varpi$ dalam
\cref{prop:b2:orbital-forcing:distance}.
\textbf{8.} Dulu $(1 - 0.001936)/(1 - 0.04395) = 0.998064/0.956049 = 1.04395$; kini $(1 -
0.000296)/(1 - 0.01686) = 0.999704/0.983139 = 1.01685$.
\textbf{9.} Dulu $1/1.04395^2 = 0.91758$; kini $1/1.01685^2 = 0.96714$.
\textbf{10.} Dulu $1 - 2 \times 0.04395 = 0.91210$; kini $1 - 2 \times 0.01686 = 0.96628$:
pendekatan itu meleset sebesar $0.0055$ dulu, $0.0009$ kini.
\textbf{11.} Dulu $1 - 0.91758 = \qty{8.2}{\percent}$ lebih lemah, kini \qty{3.3}{\percent}.
\textbf{12.} Pada solstis $\delta = \varepsilon$: \ang{22.52} dulu, \ang{23.45} kini.
\textbf{13.} Dulu $\cos h_0 = -2.14451 \times 0.41462 = -0.88916$, $h_0 = \ang{152.77} =
\qty{2.6663}{rad}$; kini $\cos h_0 = -2.14451 \times 0.43378 = -0.93023$, $h_0 = \ang{158.47} =
\qty{2.7659}{rad}$.
\textbf{14.} $2h_0/\ang{15}$ per jam: dulu $20.37$~jam, kini $21.13$~jam.
\textbf{15.} Dulu $2.6663 \times 0.90631 \times 0.38301 + 0.42262 \times 0.92375 \times 0.45759 =
0.92554 + 0.17864 = 1.10418$ (\num{1.10417} tanpa pembulatan), dan $432.90 \times 1.10417 =
\qty{478.0}{W.m^{-2}}$; kini $2.7659 \times 0.90631 \times 0.39795 + 0.42262 \times 0.91741 \times
0.36697 = 0.99757 + 0.14228 = 1.13985$ (\num{1.13982} tanpa pembulatan), dan $432.90 \times 1.13982 = \qty{493.4}{W.m^{-2}}$.
\textbf{16.} Dulu $478.0 \times 0.91758 = \qty{438.6}{W.m^{-2}}$; kini $493.4 \times 0.96714 =
\qty{477.2}{W.m^{-2}}$, \qty{38.6}{W.m^{-2}} lebih banyak daripada dulu.
\textbf{17.} Pada pertengahan Juli bujur Matahari adalah \ang{120}: $\sin\delta = \sin\varepsilon\sin\ang{120}$,
deklinasi sekitar \ang{20} alih-alih \ang{23}, sehingga siang lebih pendek dan Matahari lebih rendah: pada
\ang{65} LU insolasi harian rata-rata paling besar dekat solstis, ketika deklinasinya paling besar.
\textbf{18.} Dengan jarak masa kini, perubahan kemiringan sumbu saja mengurangi $(493.4 - 478.0) \times
0.96714 = \qty{14.9}{W.m^{-2}}$; dengan kemiringan sumbu yang lama, jarak saja mengurangi $478.0 \times
(0.96714 - 0.91758) = \qty{23.7}{W.m^{-2}}$. Keduanya berjumlah \qty{38.6}{W.m^{-2}}; jarak
berperan sekitar tiga perlima.
\textbf{19.} Setelah perubahan mendadak, lembaran es telah mencapai $1 - \mathrm e^{-t/\tau}$ darinya:
$1 - \mathrm e^{-0.5} = 0.39$ setelah \qty{5}{ka}, $1 - \mathrm e^{-1} = 0.63$ setelah \qty{10}{ka},
$1 - \mathrm e^{-2} = 0.86$ setelah \qty{20}{ka}.
\textbf{20.} $\omega\tau = 2\pi \times 10/23 = 2.732$: lembaran es mengikuti $1/\sqrt{1 + 2.732^2} =
0.34$ dari gaya paksa, dengan keterlambatan $\arctan(2.732)/\omega = 4.5$~ka.
\textbf{21.} Setengah siklus presesi, sekitar 11 sampai \qty{12}{ka}, kira-kira sama dengan satu waktu tanggap: sebelum
lembaran es menempuh dua pertiga jalannya, insolasi berbalik dan menariknya kembali.
\textbf{22.} Dalam model linear, es tumbuh dan menyusut dengan laju yang sama untuk gaya paksa yang
sama besarnya: gaya paksa yang naik dan yang turun memberikan tanggapan yang saling mencerminkan. Penumpukan lambat yang diikuti keruntuhan
cepat memerlukan tanggapan yang bergantung pada keadaan atau arah perubahan, yaitu tanggapan yang
tak linear.
\textbf{23.} Karbon dioksida naik sebesar $3.93\log_2(266/194) = 1.8$ sampai $3.93\log_2(282/188) =
\qty{2.3}{W.m^{-2}}$ dari glasial ke \omterm{def:b2:orbital-forcing:glacials}{interglasial}: gaya paksa di seluruh bola Bumi dan sepanjang
tahun. Karena dilepaskan ketika sirkulasi Atlantik melemah, karbon dioksida membawa perubahan musim panas
utara ke seluruh planet, dan, bersama albedo es yang menyusut, memperkuatnya.
\textbf{24.} Ketika \omterm{def:b2:orbital-forcing:glacials}{interglasial} terakhir berakhir, insolasi harian rata-rata solstis Juni pada
\ang{65} LU sekitar \textbf{\qty{439}{W.m^{-2}}}, \qty{39}{W.m^{-2}} lebih kecil daripada masa kini dan merupakan nilai
terendahnya dalam siklus itu.
\end{solution}
